\section{Graphs of Rational Functions}

\begin{procedure}[Steps for Consutructing a Sign Diagram for a Rational Function]
    Suppose $r$ is a rational function.
    \begin{enumerate}
        \item Place any values excluded from the domain $r$ on the number line with distinguishable sign.
        \item Find the zeros of $r$ and place them on the number line with the number 0 above them.
        \item Choose a test value in each of the intervals determined in the previous steps.
        \item Determine the sign of $r(x)$ for each test value in the previous step, and write that sign above the corresponding interval.
    \end{enumerate}
\end{procedure}

\begin{procedure}[Steps for Graphing Rational Functions]
    Suppose that $r$ is a rational function.
    \begin{enumerate}
        \item Find the domain of $r$.
        \item Reduce $r(x)$ to lowest terms, if applicable.
        \item Find the $x$- and $y$-intercepts of the graph of $y=r(x)$, if they exist.
        \item Determine the location of any verical asymptotes or holes in the graph, if they exist. Analyze the behavior of $r$ on either side of the vertical asymptotes, if applicable.
        \item Analyze the end behavior of $r$. Find the horizontal or slant asymptote, if one exists.
        \item Use a sign diagram and plot additional points, as needed, to sketch the graph of $y=r(x)$
    \end{enumerate}
\end{procedure}

\subsection{Exercises}
\includepdf[pages=-]{Exercises/04/02.pdf}